\makeatletter
\@ifundefined{keepfig}{\newenvironment{keepfig}{\par\edef\keepfig@pi{\the\parindent}\noindent
  \begin{minipage}{\linewidth}\parindent=\keepfig@pi\relax}{\end{minipage}\par}}{}
\makeatother
\section{Конечные поля}

Пусть $\F_p=\Z/p\Z$. Каждое поле $E$ характеристики $p$ содержит ровно одну копию $\F_p$~--- своё простое подполе
$\{m\cdot1_E\mid m\in\Z\}$ (\rf{1.4}); мы отождествляем его с $\F_p$. В частности, многочлен $X^q-X$ при $q=p^n$ имеет
коэффициенты в $\F_p$, над каким бы полем характеристики $p$ мы его ни рассматривали.

Пусть $E$~--- поле степени $n$ над $\F_p$. Как векторное пространство размерности $n$ над $\F_p$, оно состоит из
$q=p^n$ элементов, и $E^\times$~--- группа порядка $q-1$. По теореме Лагранжа $a^{q-1}=1$ для $a\ne0$, т.\,е.\ все $q$
элементов $E$ (включая $0$)~--- корни многочлена $X^q-X$ степени $q$. Значит, $X^q-X=\prod_{a\in E}(X-a)$, и $E$~--- поле
разложения $X^q-X$ над $\F_p$; любые два поля из $q$ элементов изоморфны (\rf{2.13}). Обратно, поле разложения $X^q-X$
над $\F_p$ состоит ровно из $q$ корней этого многочлена (упражнение 2-4: корни различны, так как производная равна
$-1$, и образуют подполе, так как $x\mapsto x^q$~--- гомоморфизм колец; в частности, $(a-b)^q=a^q+(-b)^q=a^q-b^q$:
при нечётном $p$ число $q$ нечётно, а при $p=2$ имеем $-b=b$). Мы будем свободно пользоваться этими фактами.

\begin{proposition}[4.19]
Каждое расширение конечных полей простое.
\end{proposition}

\begin{proof}
Пусть $F\subset E$~--- конечные поля. Группа $E^\times$~--- конечная подгруппа мультипликативной группы поля, поэтому
она циклическая (упражнение 1-3(e)). Если $\zeta$~--- её образующая, то каждый ненулевой элемент $E$~--- степень
$\zeta$, и $E=F[\zeta]$.
\end{proof}

\begin{remark}[4.19a]
(1) Конечность $E$ существенна: бесконечное алгебраическое расширение, например $\Fpb$ над $\F_p$ (\rf{4.24}), не может быть
простым, так как $F[\alpha]$ для алгебраического $\alpha$ конечно над $F$ (\rf{1.27c}).

(2) Образующая $E$ над $F$ не обязана порождать группу $E^\times$. Так, $\F_9=\F_3[i]$, $i^2=-1$ (многочлен $X^2+1$ не
имеет корней в $\F_3$), но $i$ имеет порядок $4$, а $|\F_9^\times|=8$. Образующая $\F_9^\times$~--- например, $1+i$:
$(1+i)^2=2i$, $(1+i)^4=(2i)^2=-4=-1$, так что её порядок равен $8$.
\end{remark}

\begin{proposition}[4.20]
Для каждой степени $q=p^n$ простого числа $p$ существует поле $\F_q$ из $q$ элементов. Каждое такое поле~--- поле
разложения $X^q-X$ над $\F_p$, и потому любые два из них изоморфны. Расширение $\F_q/\F_p$~--- расширение Галуа, и его
группа Галуа~--- циклическая группа порядка $n$, порождённая \emph{автоморфизмом Фробениуса} $\sigma(a)=a^p$.
\end{proposition}

\begin{proof}
Первые два утверждения доказаны выше. Поле $\F_q$~--- поле разложения сепарабельного многочлена $X^q-X$, поэтому
$\F_q/\F_p$~--- расширение Галуа (\rf{3.10}) степени $n$, и $|\Gal(\F_q/\F_p)|=n$. Отображение $\sigma\colon a\mapsto a^p$~---
эндоморфизм Фробениуса; оно инъективно и, так как $\F_q$ конечно, биективно, т.\,е.\ является автоморфизмом (\rf{1.4}), и оно
тождественно на $\F_p$ (малая теорема Ферма). Найдём его порядок. Имеем $\sigma^n(a)=a^{p^n}=a^q=a$, т.\,е.\ $\sigma^n=\id$. Если
$\sigma^m=\id$ при $1\le m<n$, то все $q=p^n$ элементов $\F_q$~--- корни многочлена $X^{p^m}-X$, у которого не больше
$p^m<p^n$ корней,~--- противоречие. Значит, $\sigma$ имеет порядок $n$ и порождает группу порядка $n$, т.\,е.\
$\Gal(\F_q/\F_p)=\langle\sigma\rangle$. (Иначе: неподвижное поле $\langle\sigma\rangle$ состоит из корней $X^p-X$, т.\,е.\ равно
$\F_p$, и по теореме Артина \rf{3.11} $\Gal(\F_q/\F_p)=\langle\sigma\rangle$.)
\end{proof}

Таким образом, внутри поля $\F_{p^n}$ подполе $\F_p$~--- это множество неподвижных точек Фробениуса $\sigma$, а всё $\F_{p^n}$~---
множество неподвижных точек его $n$-й итерации $\sigma^n$. Это описание переносится на любое поле, содержащее
$\F_{p^n}$ (см.\ \rf{4.23}).

\begin{corollary}[4.21]
Пусть $E$~--- поле из $p^n$ элементов. Для каждого натурального делителя $m$ числа $n$ поле $E$ содержит ровно одно
подполе из $p^m$ элементов, а других подполей у $E$ нет.
\end{corollary}

\begin{proof}
По \rf{4.20} $E/\F_p$~--- расширение Галуа с циклической группой $\langle\sigma\rangle$ порядка $n$. В циклической группе порядка
$n$ для каждого делителя $d$ числа $n$ есть ровно одна подгруппа порядка $d$, других подгрупп нет. Для $m\mid n$
подгруппа порядка $n/m$ (т.\,е.\ индекса $m$)~--- это $\langle\sigma^m\rangle$. По основной теореме (\rf{3.17}) подполя $E$
взаимно однозначно соответствуют подгруппам, и подполе $E^{\langle\sigma^m\rangle}$ имеет степень $(\langle\sigma\rangle:\langle\sigma^m\rangle)=m$ над
$\F_p$, т.\,е.\ состоит из $p^m$ элементов. Любое подполе $K\subset E$ имеет некоторую степень $m$ над $\F_p$, и
$m\mid n$ по \rf{1.20}; значит, $K=E^{\langle\sigma^m\rangle}$.
\end{proof}

Подполе $E^{\langle\sigma^m\rangle}$ описывается явно: это $\{a\in E\mid a^{p^m}=a\}$, корни многочлена $X^{p^m}-X$.

\begin{center}
\begin{keepfig}
\centering
\begin{tikzpicture}[every node/.style={font=\small}]
\begin{scope}
\node (a1) at (0,0) {$\F_p$}; \node (a2) at (-1.3,1.2) {$\F_{p^2}$}; \node (a3) at (1.3,1.2) {$\F_{p^3}$};
\node (a4) at (-2.6,2.4) {$\F_{p^4}$}; \node (a6) at (0,2.4) {$\F_{p^6}$}; \node (a12) at (-1.3,3.6) {$\F_{p^{12}}$};
\draw (a1)--(a2)--(a4)--(a12); \draw (a1)--(a3)--(a6)--(a12); \draw (a2)--(a6);
\end{scope}
\begin{scope}[xshift=7cm]
\node (b1) at (0,3.6) {$\langle\sigma\rangle$}; \node (b2) at (-1.3,2.4) {$\langle\sigma^2\rangle$}; \node (b3) at (1.3,2.4) {$\langle\sigma^3\rangle$};
\node (b4) at (-2.6,1.2) {$\langle\sigma^4\rangle$}; \node (b6) at (0,1.2) {$\langle\sigma^6\rangle$}; \node (b12) at (-1.3,0) {$\{1\}$};
\draw (b1)--(b2)--(b4)--(b12); \draw (b1)--(b3)--(b6)--(b12); \draw (b2)--(b6);
\end{scope}
\draw[cnew,<->] (1.9,1.8) -- node[above,font=\footnotesize] {$K\mapsto\Gal(\F_{p^{12}}/K)$} (4.2,1.8);
\end{tikzpicture}
\figcap{Подполя $\F_{p^{12}}$ (слева) и подгруппы $\Gal(\F_{p^{12}}/\F_p)=\langle\sigma\rangle\cong\Z/12\Z$ (справа): подполе из $p^m$
элементов отвечает подгруппе $\langle\sigma^m\rangle$ порядка $12/m$. Решётка подполей~--- решётка делителей числа $12$; например,
$\F_{p^4}\not\subset\F_{p^6}$.}
\end{keepfig}
\end{center}

\begin{corollary}[4.22]
Пусть $f\in\F_p[X]$~--- унитарный неприводимый многочлен степени $d$.
\begin{enumerate}[label=\textup{(\alph*)}]
\item $f$ делит $X^{p^n}-X$ тогда и только тогда, когда $d\mid n$; в этом случае $f$ входит в разложение $X^{p^n}-X$ на
неприводимые множители ровно один раз.
\item Если $\alpha$~--- корень $f$ в каком-либо расширении $\F_p$, то $\F_p[\alpha]$~--- поле разложения $f$, его степень над
$\F_p$ равна $d$, и корни $f$~--- это $d$ различных элементов $\alpha,\alpha^p,\alpha^{p^2},\dots,\alpha^{p^{d-1}}$.
\end{enumerate}
\end{corollary}

\begin{proof}
Пусть $\alpha$~--- корень $f$ в некотором расширении (например, класс $x$ в $\F_p[X]/(f)$). Поле $K=\F_p[\alpha]$ имеет
степень $d$ над $\F_p$ (\rf{1.27c}), т.\,е.\ состоит из $p^d$ элементов.

(b) Все элементы $K$~--- корни $X^{p^d}-X$, поэтому $\alpha^{p^d}=\alpha$, и $f$, минимальный многочлен $\alpha$, делит
$X^{p^d}-X=\prod_{a\in K}(X-a)$. Значит, $f$ раскладывается в $K$ на линейные множители, и $K$, порождённое одним корнем,~---
поле разложения $f$ степени ровно $d$. Автоморфизм Фробениуса $\sigma$ поля $K$ переводит корни $f$ в корни ($f$ имеет
коэффициенты в $\F_p$), так что $\sigma^i(\alpha)=\alpha^{p^i}$~--- корни $f$. При $0\le i<j<d$ они различны: иначе $\sigma^{j-i}$
оставлял бы на месте $\alpha$, а с ним и всё $K=\F_p[\alpha]$, но $\sigma$ имеет порядок $d$ (\rf{4.20}). Это $d$ различных
корней многочлена степени $d$, т.\,е.\ все его корни.

(a) Если $d\mid n$, то, применяя равенство $\alpha^{p^d}=\alpha$ $n/d$ раз, получаем $\alpha^{p^n}=\alpha$; значит, $f$ делит
$X^{p^n}-X$ (\rf{1.27c}). Обратно, пусть $f\mid X^{p^n}-X$. Тогда $\alpha^{p^n}=\alpha$, т.\,е.\ $\sigma^n(\alpha)=\alpha$, и $\sigma^n$
тождественен на $K=\F_p[\alpha]$; так как $\sigma$ имеет порядок $d$, имеем $d\mid n$. Наконец, $X^{p^n}-X$ сепарабелен (его
производная равна $-1$, \rf{2.21}), поэтому никакой неприводимый многочлен не входит в его разложение дважды: если
$h=f^2g$, то $h'=f\cdot(2f'g+fg')$, и $f$ был бы общим делителем $h$ и $h'$.
\end{proof}

В (b) степень поля разложения равна $d$ в точности, а не просто не превосходит $d$: поле, порождённое \emph{одним}
корнем неприводимого многочлена над конечным полем, уже содержит все остальные. Над $\Q$ это, конечно, не так
($\Q[\sqrt[3]2]$ не содержит корней $\omega\sqrt[3]2$, $\omega^2\sqrt[3]2$).

\begin{aside}[4.22a: почему нет канонического изоморфизма]
Пусть $d\mid n$ и $f$~--- неприводимый многочлен степени $d$. Поле $\F_p[X]/(f)$ из $p^d$ элементов изоморфно
\emph{единственному} подполю $\F_{p^d}\subset\F_{p^n}$ (\rf{4.21}): образ любого $\F_p$-гомоморфизма $\F_p[X]/(f)\to\F_{p^n}$
состоит из $p^d$ элементов и потому совпадает с $\F_{p^d}$. Сам же гомоморфизм определяется образом класса $x$, а им
может быть любой корень $f$ в $\F_{p^n}$, т.\,е.\ любой из $d$ элементов $\alpha,\alpha^p,\dots,\alpha^{p^{d-1}}$ (\rf{2.1}(b), \rf{4.22}(b)). Получаются
$d$ различных изоморфизмов $\F_p[X]/(f)\to\F_{p^d}$, отличающихся друг от друга степенями Фробениуса~--- автоморфизмами
$\F_{p^d}$, и ни один из них ничем не выделен: корни $f$ равноправны (их переставляет группа Галуа). Так и над $\Q$
есть два изоморфизма $\Q[X]/(X^2+1)\to\Q[i]$, $x\mapsto\pm i$. Канонично лишь \emph{подполе} $\F_{p^d}$, а не его
отождествление с $\F_p[X]/(f)$; по той же причине разные неприводимые многочлены степени $d$ дают поля $\F_p[X]/(f)$,
изоморфные, но не канонически.
\end{aside}

\begin{center}
\begin{keepfig}
\centering
\begin{tikzpicture}[every node/.style={font=\small}]
\node[cold] (Q) at (0,0) {$\F_p[X]/(f)$};
\node (Fd) at (5.6,0) {$\F_{p^d}$}; \node (Fn) at (8,0) {$\F_{p^n}$};
\draw[cnew,->] (Q.east) to[bend left=32] node[above,font=\footnotesize] {$x\mapsto\alpha$} (Fd.west);
\draw[cnew,->] (Q.east) -- node[above,font=\footnotesize] {$x\mapsto\alpha^p$} (Fd.west);
\draw[cnew,->] (Q.east) to[bend right=32] node[below,font=\footnotesize] {$x\mapsto\alpha^{p^{d-1}}$} (Fd.west);
\node[font=\footnotesize] at (2.9,-0.32) {$\vdots$};
\draw[right hook->] (Fd) -- (Fn);
\node[font=\footnotesize,align=center] at (6.8,-1.1) {единственное подполе\\ из $p^d$ элементов};
\end{tikzpicture}
\figcap{Образ любого вложения $\F_p[X]/(f)\to\F_{p^n}$~--- одно и то же подполе $\F_{p^d}$, но самих вложений $d$, по числу корней
$f$; они отличаются автоморфизмами Фробениуса поля $\F_{p^d}$.}
\end{keepfig}
\end{center}

\begin{theorem}[4.22b]
В $\F_p[X]$
\[
X^{p^n}-X=\prod_{d\mid n}\ \prod_{\substack{f\ \text{унитарный неприводимый},\\ \deg f=d}}f .
\]
Иначе говоря, разложив $X^{p^n}-X$ на неприводимые множители, мы получим по одному разу все унитарные неприводимые
многочлены, степени которых делят $n$.
\end{theorem}

\begin{proof}
По \rf{4.22}(a) каждый унитарный неприводимый $f$ степени $d\mid n$ делит $X^{p^n}-X$; различные унитарные неприводимые
многочлены попарно взаимно просты, поэтому их произведение делит $X^{p^n}-X$. Обратно, каждый неприводимый унитарный
множитель $X^{p^n}-X$ имеет по \rf{4.22}(a) степень, делящую $n$, и входит в разложение один раз. Значит, разложение
$X^{p^n}-X$ (унитарного) на неприводимые множители~--- в точности это произведение.
\end{proof}

Сравнение степеней даёт формулу для числа $N_p(d)$ унитарных неприводимых многочленов степени $d$ над $\F_p$:
\begin{equation}\label{eq:c4-gauss}
p^n=\sum_{d\mid n}d\,N_p(d).
\end{equation}

\begin{corollary}[4.22c]
\textup{(a)} Если $\ell$ простое, то $N_p(\ell)=\dfrac{p^\ell-p}\ell$; в частности, $\ell\mid p^\ell-p$ \textup(малая теорема
Ферма\textup).

\textup{(b)} \textup(Гаусс\textup) $N_p(n)=\dfrac1n\sum_{d\mid n}\mu(d)\,p^{n/d}$, где $\mu$~--- функция Мёбиуса, и $N_p(n)\ge1$ для
всех $n\ge1$.
\end{corollary}

\begin{proof}
(a) При $n=\ell$ делители~--- $1$ и $\ell$, и \eqref{eq:c4-gauss} даёт $p^\ell=N_p(1)+\ell N_p(\ell)=p+\ell N_p(\ell)$.
(b) Формула получается из \eqref{eq:c4-gauss} обращением Мёбиуса (функции $n\mapsto p^n$ и $d\mapsto dN_p(d)$ связаны
суммированием по делителям). Неравенство: по \eqref{eq:c4-gauss}
$nN_p(n)=p^n-\sum_{d\mid n,\,d<n}dN_p(d)\ge p^n-\sum_{d\le n/2}p^d>p^n-p^{n/2+1}\ge0$ при $n\ge2$ (так как $dN_p(d)\le p^d$,
а $\sum_{d\le n/2}p^d<p^{\lfloor n/2\rfloor+1}/(p-1)\le p^{n/2+1}$); при $n=1$ $N_p(1)=p$. Это ещё одно доказательство
существования поля из $p^n$ элементов: $\F_p[X]/(f)$ при $\deg f=n$.
\end{proof}

\begin{example}[4.22d]
Над $\F_2$: $N_2(1),\dots,N_2(8)=2,1,2,3,6,9,18,30$, и
\begin{align*}
X^8-X&=X(X+1)(X^3+X+1)(X^3+X^2+1),\\
X^{16}-X&=X(X+1)(X^2+X+1)(X^4+X+1)(X^4+X^3+1)(X^4+X^3+X^2+X+1).
\end{align*}
Во втором разложении нет кубических множителей: $3\nmid4$, и $\F_8\not\subset\F_{16}$; зато есть $X^2+X+1$, корни
которого вместе с $0,1$ образуют $\F_4\subset\F_{16}$. Над $\F_3$:
$X^9-X=X(X+1)(X+2)(X^2+1)(X^2+X+2)(X^2+2X+2)$, т.\,е.\ $N_3(2)=3$. Над $\F_5$ разложение $X^{625}-X$ (так как $625=5^4$)
содержит все $5$ унитарных многочленов степени $1$, все $N_5(2)=10$ неприводимых степени $2$ и все $N_5(4)=150$
неприводимых степени $4$ ($5+2\cdot10+4\cdot150=625$). Такие разложения быстро вычисляются на компьютере (в PARI~---
командой \texttt{factormod(x\^{}625-x,5)}), и это простой способ получить список всех неприводимых многочленов данной
степени.
\end{example}

\begin{proposition}[4.23]
Пусть $\Fpb$~--- алгебраическое замыкание $\F_p$. Для каждого $n\ge1$ поле $\Fpb$ содержит ровно одно подполе $\F_{p^n}$ из
$p^n$ элементов, и оно состоит из корней $X^{p^n}-X$. При этом
\[
\F_{p^m}\subset\F_{p^n}\iff m\mid n,
\]
и в этом случае $\F_{p^n}/\F_{p^m}$~--- расширение Галуа, группа которого циклическая порядка $n/m$ и порождена
автоморфизмом $x\mapsto x^{p^m}$.
\end{proposition}

\begin{proof}
Многочлен $X^{p^n}-X$ сепарабелен и раскладывается в $\Fpb$ на линейные множители, так что у него $p^n$ различных корней в
$\Fpb$; они образуют подполе (как в упражнении 2-4). Если $K\subset\Fpb$~--- любое подполе из $p^n$ элементов, то
$a^{p^n}=a$ для всех $a\in K$ (теорема Лагранжа для $K^\times$), т.\,е.\ $K$ содержится в множестве корней и, состоя из
стольких же элементов, совпадает с ним.

Если $\F_{p^m}\subset\F_{p^n}$ и $[\F_{p^n}:\F_{p^m}]=e$, то $p^n=(p^m)^e$, т.\,е.\ $n=me$ и $m\mid n$. Обратно, если $m\mid n$, то
каждый корень $X^{p^m}-X$ является корнем $X^{p^n}-X$ (применяем $a^{p^m}=a$ $n/m$ раз), т.\,е.\ $\F_{p^m}\subset\F_{p^n}$.
Пусть $\tau(x)=x^{p^m}$~--- автоморфизм $\F_{p^n}$ ($m$-я степень Фробениуса $\sigma$). Его неподвижные точки~--- корни $X^{p^m}-X$, т.\,е.\
$\F_{p^m}$; по теореме Артина (\rf{3.11}) $\F_{p^n}/\F_{p^m}$~--- расширение Галуа с группой $\langle\tau\rangle$, а её порядок равен
$[\F_{p^n}:\F_{p^m}]=n/m$.
\end{proof}

Итак, частично упорядоченное по включению множество конечных подполей $\Fpb$ изоморфно множеству натуральных чисел,
упорядоченному по делимости, а $\F_{p^n}=\{a\in\Fpb\mid\sigma^n(a)=a\}$, где $\sigma\colon a\mapsto a^p$~--- автоморфизм Фробениуса
поля $\Fpb$ (\rf{1.4}). Всё $\Fpb$~--- объединение своих конечных подполей (каждый элемент алгебраичен над $\F_p$ и порождает
конечное поле $\F_p[a]$).

\begin{proposition}[4.24]
У поля $\F_p$ есть алгебраическое замыкание.
\end{proposition}

\begin{proof}
Выберем последовательность натуральных чисел $1=n_1<n_2<n_3<\dots$, в которой $n_i\mid n_{i+1}$ и каждое натуральное
число делит некоторое $n_i$; например, $n_i=i!$. Построим поля индуктивно: $F_1=\F_p$, а $F_i$~--- поле разложения
многочлена $X^{p^{n_i}}-X$ над $F_{i-1}$.

\emph{$F_i$ состоит из $p^{n_i}$ элементов.} Множество $S$ корней $X^{p^{n_i}}-X$ в $F_i$~--- подполе из $p^{n_i}$ элементов
(упражнение 2-4). По индукции $F_{i-1}$ состоит из $p^{n_{i-1}}$ элементов, т.\,е.\ $a^{p^{n_{i-1}}}=a$ для $a\in F_{i-1}$, а так
как $n_{i-1}\mid n_i$, то и $a^{p^{n_i}}=a$: $F_{i-1}\subset S$. Поле $F_i$ порождено над $F_{i-1}$ корнями, т.\,е.\ элементами
$S$; значит, $F_i=S$.

Получаем цепочку $\F_p=F_1\subset F_2\subset F_3\subset\dots$ и полагаем $\F=\bigcup_iF_i$ (см.\ \rf{4.25}). Объединение цепочки полей~--- поле
(любые два элемента лежат в одном $F_i$), и оно алгебраично над $\F_p$ (каждый элемент лежит в конечном поле $F_i$).
Пусть $f\in\F_p[X]$ и $g$~--- неприводимый унитарный множитель $f$ степени $d$. Тогда $d\mid n_i$ для некоторого $i$, по
\rf{4.22}(a) $g$ делит $X^{p^{n_i}}-X$, а этот многочлен раскладывается на линейные множители в $F_i$. Значит, каждый
многочлен из $\F_p[X]$ раскладывается в $\F$ на линейные множители, и по \rf{1.45} $\F$~--- алгебраическое замыкание $\F_p$.
\end{proof}

\begin{remark}[4.25]
(1) Поля $F_i$ в \rf{4.24} не являются подмножествами одного заранее данного множества: <<включение>> $F_{i-1}\subset F_i$~---
это вложение, полученное при построении поля разложения. Поэтому объединение требует пояснения. Можно сослаться на
аксиому объединения теории множеств или поступить наивно: взять дизъюнктное объединение $S=\bigsqcup F_i$ и считать
$a\in F_i$ и $b\in F_j$, $i\le j$, эквивалентными, если $b$~--- образ $a$ при вложении $F_i\to F_j$; тогда $\F=S/{\sim}$, и
операции в $\F$ определяются в достаточно больших $F_j$. (На языке категорий $\F$~--- копредел, или индуктивный
предел, цепочки $F_1\to F_2\to\dots$.)

(2) Любые два алгебраических замыкания $\F_p$ изоморфны (\rf{6.6}), но не канонически: например, изоморфизм можно
изменить на любой автоморфизм, и группа $\Aut(\Fpb/\F_p)$ огромна (она изоморфна $\Zh$ и несчётна, \rf{7.17}, \rf{7.26b}). Точно так же любые два поля
из $q$ элементов изоморфны, но, вообще говоря, не канонически (\rf{4.22a}). Однако после того как алгебраическое замыкание
$\Fpb$ выбрано, в нём есть \emph{единственное} подполе из $q$ элементов (\rf{4.23}), и обозначение $\F_q$ становится
однозначным.

(3) Конечные поля открыл Галуа (статья <<О теории чисел>>, 1830): он ввёл поля из $p^n$ элементов, присоединяя к $\F_p$
<<мнимый>> корень неприводимого сравнения, и установил их основные свойства, например существование образующей
мультипликативной группы. Поэтому конечные поля называют также \emph{полями Галуа} и обозначают $\mathrm{GF}(q)$.
\end{remark}

\subsection*{Группа Галуа многочлена над конечным полем}

Всё сказанное с заменой $p$ на $q=p^k$ переносится на расширения $\F_q\subset\F_{q^n}$ (подполя $\Fpb$): по \rf{4.23} (с
$p^m=q$) расширение $\F_{q^n}/\F_q$~--- расширение Галуа с циклической группой порядка $n$, порождённой
\emph{автоморфизмом Фробениуса над $\F_q$} $\varphi_q\colon x\mapsto x^q$. Отсюда~--- полное описание групп Галуа
многочленов над конечными полями; на нём основана теорема Дедекинда (\rf{4.28}).

Напомним, что \emph{цикловой тип} перестановки конечного множества~--- набор длин её независимых циклов (с учётом
циклов длины $1$; ср.\ таблицу после \rf{4.18a}).
Циклы перестановки $\pi$~--- это в точности орбиты циклической группы $\langle\pi\rangle$: орбита элемента $x$ есть
$\{x,\pi x,\pi^2x,\dots\}$, и $\pi$ переставляет её элементы по кругу.

\begin{proposition}[4.25a]
Пусть $f\in\F_q[X]$~--- сепарабельный многочлен степени $m$ и $f=f_1\cdots f_r$~--- его разложение на неприводимые
множители, $m_i=\deg f_i$. Тогда группа Галуа $G_f$ циклическая и порождена автоморфизмом Фробениуса $\varphi_q\colon x\mapsto x^q$;
как перестановка корней $f$, $\varphi_q$ есть произведение $r$ независимых циклов длин $m_1,\dots,m_r$, т.\,е.\ имеет
цикловой тип $(m_1,\dots,m_r)$: корни каждого $f_i$ образуют один цикл $\beta\mapsto\beta^q\mapsto\beta^{q^2}\mapsto\dots$ длины $m_i$.
Порядок $G_f$, т.\,е.\ степень поля разложения $f$ над $\F_q$, равен наименьшему общему кратному чисел $m_1,\dots,m_r$.
Каждый неприводимый многочлен над $\F_q$ сепарабелен, так что сепарабельность $f$ означает лишь, что среди $f_i$ нет
ассоциированных \textup(отличающихся постоянным множителем\textup).
\end{proposition}

\begin{proof}
Пусть $E$~--- поле разложения $f$ над $\F_q$ и $N=[E:\F_q]$, так что $|E|=q^N$. Отображение $\varphi_q$~--- степень
эндоморфизма Фробениуса, т.\,е.\ инъективный гомоморфизм конечного поля $E$ в себя и потому его автоморфизм (\rf{1.4});
он тождествен на $\F_q$ ($a^q=a$ для всех $a\in\F_q$, теорема Лагранжа для $\F_q^\times$). Его неподвижные точки~--- корни
$X^q-X$ в $E$; их не больше $q$, и все $q$ элементов $\F_q$~--- такие корни. Значит, $E^{\langle\varphi_q\rangle}=\F_q$, и по
теореме Артина (\rf{3.11}) $G_f=\Gal(E/\F_q)=\langle\varphi_q\rangle$ (если $E\subset\Fpb$, то это \rf{4.23} с $p^m=q$: $E=\F_{q^N}$).

По \rf{3.27d} орбиты $G_f$ на корнях~--- множества корней множителей $f_i$. Орбита циклической группы $\langle\varphi_q\rangle$,
содержащая корень $\beta$ множителя $f_i$, есть $\{\beta,\beta^q,\beta^{q^2},\dots\}$, т.\,е.\ один цикл перестановки $\varphi_q$; его
длина равна размеру орбиты, т.\,е.\ числу корней $f_i$, а оно равно $m_i$, так как $f$ сепарабелен. Порядок
перестановки~--- наименьшее общее кратное длин её независимых циклов, а $G_f$ действует на корнях точно.

Наконец, пусть $g\in\F_q[X]$ неприводим степени $d$ и $\beta$~--- его корень в некотором расширении. Поле $\F_q[\beta]$ состоит
из $q^d$ элементов, поэтому $\beta^{q^d}=\beta$, и $g$ (с точностью до постоянного множителя~--- минимальный многочлен $\beta$)
делит $X^{q^d}-X$. Этот многочлен сепарабелен (его производная равна $-1$), значит, сепарабелен и $g$ (при $q=p$ это
\rf{4.22}(b)). Неассоциированные неприводимые многочлены общих корней не имеют (каждый из них с точностью до постоянного
множителя~--- минимальный многочлен любого своего корня, \rf{1.27c}). Поэтому кратный корень у $f$ бывает, только если
некоторый неприводимый множитель входит в разложение $f$ дважды.
\end{proof}

Например, по модулю $2$ имеем $X^5-X-1\equiv(X^2+X+1)(X^3+X^2+1)$ (оба множителя неприводимы: нет корней в $\F_2$). Значит,
над $\F_2$ Фробениус переставляет корни как произведение цикла длины $2$ и цикла длины $3$, группа Галуа~--- $\Z/6\Z$,
поле разложения~--- $\F_{64}$. Теорема Дедекинда утверждает, что перестановка с таким же цикловым типом есть и в группе
Галуа многочлена $X^5-X-1$ над $\Q$ (\rf{4.30}). Для неприводимого $f$ степени $d$ предложение \rf{4.25a}~--- это \rf{4.22}(b): $\varphi_p$
переставляет корни $\alpha,\alpha^p,\dots,\alpha^{p^{d-1}}$ по кругу.
